\section{Experimental settings and metrics}\label{app:settings}
Training and functional evaluations used a cluster of NVIDIA A100 GPUs. Each task used a single GPU, with at most 30 GPUs in use simultaneously. The post-hoc hyperspherical-energy measurements used NVIDIA L40 GPUs.

\begin{table}[!htbp]
\caption{Study overview. Each comparison uses a shared pretrained backbone. Divergent runs are reported with the corresponding results.}
\label{tab:studies}
\begin{center}
\setlength{\tabcolsep}{3pt}
\begin{tabular*}{\linewidth}{@{\extracolsep{\fill}}lll@{}}\toprule
Backbone & Adaptation & Primary retention view \\\midrule
Qwen2.5-7B & MetaMath / six methods & General-text NLL \\
Llama-3.1-8B & MetaMath / six methods & General-text NLL \\
Qwen2.5-7B & CodeFeedback / six methods & General-text NLL \\
Llama-3.1-8B & CodeFeedback / six methods & General-text NLL \\
FLUX-4B & Cat / six methods & Image drift / CLIP \\
FLUX-4B & 28 objects & Flow loss / drift / CLIP \\\bottomrule
\end{tabular*}
\end{center}
\end{table}

\subsection{Language adaptation}
Qwen2.5-7B and Llama-3.1-8B use 20,000 MetaMath examples \citep{yu2024metamath}, microbatch size two, gradient accumulation over 16 steps, sequence length 768, 625 optimizer updates, cosine decay, 3\% warmup, gradient clipping at one and zero weight decay. LoRA uses $\alpha=2r$ and zero dropout. OFT uses unshared coordinate blocks and OFTv2's Cayley--Neumann parameterization with five terms \citep{qiu2025oftv2}. The seven projection types are query, key, value, output, up, gate and down, giving 196 adapted matrices in Qwen and 224 in Llama. LRs are $\{1,3,10,30,100\}\times10^{-5}$ and seeds are 13, 37 and 73. 

All six methods use AdamW \citep{loshchilov2019adamw}. LoRA-family ranks are 7/14/28 for Qwen and 7/15/30 for Llama. OFT block sizes are 32/64/128, and HRA uses 16/32/64 Householder vectors. For each model/method/capacity, mean accuracy on the fixed 1,000-item GSM8K development panel selects one LR from fully completed configurations, with ties favoring the lower LR. The resulting 108 checkpoints are evaluated on the official 1,319-item test set. The diverged Qwen PiSSA configuration is excluded from LR selection and complete-configuration frontiers. Appendix~\ref{app:language} reports all selected outcomes and the later HRA LR extension separately.

Coding uses the same two backbones, six methods and five-rate grid at one capacity per method. Its held-out response-NLL selection, completion-based testing and full settings are specified in Appendix~\ref{app:coding}.

Standard and intervention retention use FineWiki-200 \citep{penedo2025finewiki}, with contexts truncated to 768 tokens. The intervention panel evaluates GSM8K development accuracy on 1,000 items. General NLL and full-vocabulary KL use all 200 FineWiki documents, with token-pooled aggregation. General NLL scores valid next-token positions. All intervention comparisons use the same evaluation inputs and checkpoint-local references.

\subsection{Diffusion personalization}
FLUX.2-klein-base-4B uses AdamW, batch size two, weight decay $10^{-4}$, gradient clipping at one, constant LR and 750 updates. The benchmark-cat set has 14 training, two validation and four test images. Four approximately matched parameter pairs are r4/b32, r8/b64, r16/b128 and r32/b256. LoRA uses $\alpha=r$. OFT uses unshared coordinate blocks with OFTv2's Cayley--Neumann parameterization \citep{qiu2025oftv2}. The LR grid is $\{0.5,1,3,5,10,30,50,100\}\times10^{-5}$, with seeds 0, 1 and 2. Generation uses resolution 512, 20 inference steps and guidance 3.5. For each method/capacity, we select a shared LR by the mean validation DINO at step 700, the last validation evaluation, then report test and retention scores at final step 750. The two validation and four test images are disjoint.

The complete six-method cat comparison uses LoRA/DoRA/PiSSA/MiLoRA ranks 4/8/16/32, OFT blocks 32/64/128/256, and HRA 16/32/62/124 Householder vectors. All 576 runs across six methods, four capacities and eight LRs are complete. 

Target DINOv2 similarity averages the paired generated/reference test-image cosines over four examples and three generation repeats. General and class drift compare paired base/adapted generations using $(1-\mathrm{cosine})/2$, over five prompts in each bank. General CLIP is normalized text/image cosine averaged over five general prompts. Training and generation randomness contribute to the reported seed variation.

The 28-object study uses three training images, one validation image and one or two held-out test images per object. It crosses the same four capacity pairs with LRs $\{3,5,10\}\times10^{-5}$ and seeds 0, 1 and 2, for 2,016 runs. For each method/capacity, we select one LR by mean step-700 validation DINO across all 28 objects and seeds, breaking ties toward the smaller LR. That LR is shared across objects and seeds. Test DINO, averaged over three generation repeats, and retention are measured at step 750. The retention bank uses 200 COCO image/caption pairs and four scheduler-index fractions, $\{0.125,0.375,0.625,0.875\}$. Image drift and CLIP use the same four general prompts for every object, depicting an angry Mona Lisa, the Eiffel Tower at night, an astronaut riding a horse, and an avocado-shaped chair.

\subsection{Metrics and uncertainty}
\paragraph{Language likelihood.}
For document $j$ and its valid next-token positions $T_j$, general-text NLL is
\begin{equation*}
 \mathrm{NLL}(\bm\Theta)=-\frac{1}{\sum_j |T_j|}\sum_j\sum_{t\in T_j}
 \log p_{\bm\Theta}\!\left(x_{j,t}\middle|x_{j,<t}\right).
\end{equation*}
All models use identical teacher-forced contexts and masks, excluding padding. The score weights tokens equally and measures fit to the general-text bank.

\paragraph{Diffusion outputs and predictive change.}
Write $\cos(\bm a,\bm b)=\frac{\bm a^{\!\top}\bm b}{\left\|\bm a\right\|_2\left\|\bm b\right\|_2}$, with normalized DINOv2 features $\bm{\mathrm{f}}(I)$ and CLIP image/text features $\bm{\mathrm{g}}(I),\bm{\mathrm{h}}(c)$. For generated image $I_j^{\bm\Theta}$, paired base image $I_j^0$, target reference image $I_j^{\rm ref}$ and prompt $c_j$, the reported averages are
\begin{align*}
 \mathrm{DINO}&=\frac{1}{J}\sum_j\cos\!\left(\bm{\mathrm{f}}(I_j^{\bm\Theta}),\bm{\mathrm{f}}(I_j^{\rm ref})\right),
 &D_{\rm image}&=\frac{1}{J}\sum_j\frac{1-\cos\!\left(\bm{\mathrm{f}}(I_j^{\bm\Theta}),\bm{\mathrm{f}}(I_j^0)\right)}{2},\\
 \mathrm{CLIP}&=\frac{1}{J}\sum_j\cos\!\left(\bm{\mathrm{g}}(I_j^{\bm\Theta}),\bm{\mathrm{h}}(c_j)\right).
\end{align*}
Each metric uses its own stated bank of $J$ pairs. General and class drift use different prompt banks. Paired images share generation seeds. Repeats are averaged before training-seed summaries. The 28-object tables report absolute target DINO, COCO flow loss and CLIP alignment, together with base-relative image drift.

For the 28-object COCO flow loss, $\bm y_{j,k}$ contains the fixed noisy latent, scheduler time and conditioning for image/caption pair $j$ in time bucket $k$. With $d_j$ velocity coordinates, $K$ buckets, flow target $\bm u_{j,k}$ and loss weight $w_{j,k}$,
\begin{equation*}
 \mathcal L_{\rm flow}(\bm\Theta)=\frac{1}{NK}\sum_{j,k}\frac{w_{j,k}}{d_j}
 \left\|\bm v_{\bm\Theta}(\bm y_{j,k})-\bm u_{j,k}\right\|_2^2.
\end{equation*}
The rectified-flow target is noise minus clean latent \citep{liu2023rectifiedflow}. This fixed-state prediction loss is distinct from generated-image drift.

\paragraph{Uncertainty.}
For $S$ training seeds with scores $z_s$, we report $\bar z=S^{-1}\sum_s z_s$ and sample $\mathrm{SD}=\left(\sum_s(z_s-\bar z)^2/(S-1)\right)^{1/2}$. Perturbation and generation repetitions are averaged within a checkpoint first. For each 28-object method/capacity/LR setting, $z_s=28^{-1}\sum_{o=1}^{28}z_{s,o}$. The SD describes variation across training seeds.
